\section{Semantic Structure}

\todo{Understand semantics as model sets, over which distributions can be defined.}

We now investigate the \semanticStructure{} of a \firstOrderLogic{} theory, when restricting to database semantics.
This involves the representation of collections of worlds, where each has a \substitutionStructure{} as discussed above.

\subsection{World Enumerating Variables}

Given database semantics, we have a finite set of worlds, which we enumerate by tuples of variables $\catvariable$.

\begin{definition}[World enumerating variables]
    \label{def:worldEnumeratingVariables}
    Given a \firstOrderLogic{} theory with constant symbols $\worlddomain$, function symbols $\{\folfunctionof{\secatomenumerator}\wcols\secatomenumeratorin\}$ and predicate symbols $\{\folpredicateof{\atomenumerator}\wcols\atomenumeratorin\}$, we enumerate the world under database semantics by a tuple
    \begin{align*}
        \worldvariables
        = \left(\bigtimes_{\seccatenumeratorin}\bigtimes_{\indindexof{\seccatenumerator}\in[\inddim]^{\indorderof{\seccatenumerator}+1}} \catvariableof{\seccatenumerator,\indindexof{\seccatenumerator}} \right)
        \times \left(\bigtimes_{\catenumeratorin}\bigtimes_{\indindexof{\catenumerator}\in\in[\inddim]^{\indorderof{\catenumerator}}} \catvariableof{\catenumerator,\indindexof{\catenumerator}} \right)
    \end{align*}
    of boolean variables.
    In the world indexed by $\worldindices$, the grounding tensor of the function $\folfunctionof{\secatomenumerator}$ is
    \begin{align*}
        \bencodingofat{\groundingof{\folfunctionof{\secatomenumerator}}}{\indvariableof{\secatomenumerator},\indvariableof{[\inddimof{\secatomenumerator}]}}
        = \sum_{\indindexof{\seccatenumerator}\in[\inddim]^{\indorderof{\seccatenumerator}+1}\wcols\indindexof{\seccatenumerator}=1} \onehotmapofat{\indindexof{\seccatenumerator}}{\indvariableof{\seccatenumerator}}
    \end{align*}
    and of the predicate $\folpredicateof{\atomenumerator}$
    \begin{align*}
        \groundingofat{\folpredicateof{\atomenumerator}}{\indvariableof{[\inddimof{\atomenumerator}]}}
        = \sum_{\indindexof{\seccatenumerator}\in[\inddim]^{\indorderof{\atomenumerator}}\wcols\indindexof{\atomenumerator}=1}
        \onehotmapofat{\indindexof{\atomenumerator}}{\indvariableof{\atomenumerator}} \, .
    \end{align*}
\end{definition}

We further have the restriction that each basis encoded function is a directed tensor.
To restrict to those worlds, where this is true, we further introduce the validation base measure $\basemeasureat{\worldvariables}$ with coordinates
\begin{align*}
    \basemeasureat{\indexedworldvariables} =
    \begin{cases}
        1 & \ifspace \uniquantwrtof{\secatomenumeratorin}{\sum_{\indindexof{\secatomenumerator}\in[\inddim]^{\indorderof{\seccatenumerator}+1}}=1} \\
        0 & \text{else}
    \end{cases} \, .
\end{align*}

\subsection{Representation of Terms and Formulas}

Combining its \substitutionStructure{} and \semanticStructure{} we represent a \firstOrderLogic{} term $\folterm$ with arity $\inddimof{\folterm}$ as a tensor
\begin{align*}
    \bencodingofat{\folterm}{\indvariableof{\folterm},\indvariableof{[\inddimof{\folterm}]},\worldvariables} =
    \sum_{\worldindices\wcols\basemeasureat{\indexedworldvariables}=1}
    \bencodingofat{\groundingof{\folterm}}{\indvariableof{\folterm},\indvariableof{[\inddimof{\folterm}]}}
    \otimes \onehotmapofat{\worldindices}{\worldvariables} \, .
\end{align*}
and a formula as the tensor
\begin{align*}
    \folexformulawith =
    \sum_{\worldindices\wcols\basemeasureat{\indexedworldvariables}=1} \groundingofat{\folexformula}{\indvariableof{\folexformula}}
    \otimes \onehotmapofat{\worldindices}{\worldvariables} \, .
\end{align*}

From \defref{def:worldEnumeratingVariables} the representation of predicate and function symbols are directly derived from the $\worldvariables$, that is
\begin{align*}
    \folpredicateofat{\atomenumerator}{\indvariableof{[\inddimof{\atomenumerator}]},\indexedworldvariables}
    = \sum_{\indindexof{\seccatenumerator}\in[\inddim]^{\indorderof{\seccatenumerator}+1}\wcols\indindexof{\seccatenumerator}=1}
    \onehotmapofat{\indindexof{\seccatenumerator}}{\indvariableof{\seccatenumerator}}
\end{align*}
and
\begin{align*}
    \bencodingofat{\folfunctionof{\secatomenumerator}}{\indvariableof{\secatomenumerator},\indvariableof{[\inddimof{\secatomenumerator}]},\indexedworldvariables}
    = \sum_{\indindexof{\seccatenumerator}\in[\inddim]^{\indorderof{\seccatenumerator}+1}\wcols\indindexof{\seccatenumerator}=1} \onehotmapofat{\indindexof{\seccatenumerator}}{\indvariableof{\seccatenumerator}}
\end{align*}

\subsection{Case of Propositional Logic}

\PropositionalLogic{} as discussed in (see \charef{cha:logicalRepresentation}) is the special case of \firstOrderLogic{} with empty sets of constant and function symbols, and $\indorderof{\atomenumerator}=0$ for the predicate symbols.
With the convention $[]^{0}=[2]$ the worlds are enumerated by the tuple
\begin{align*}
    \worldvariables = \bigtimes_{\atomenumeratorin} \catvariableof{\atomenumerator}
\end{align*}
which we have in previous chapters abbreviated by $\shortcatindices$.


To represent logical formulas as sets of possible worlds, and distributions of worlds, we applied in \parref{par:one} one-hot encodings of possible worlds.
For the case of \propositionalLogic{}, this is
\begin{align*}
    \onehotmapofat{\worldindices}{\shortcatvariables}
    = \bigotimes_{\atomenumeratorin} \onehotmapofat{\catindexof{\atomenumerator}}{\catvariableof{\atomenumerator}} \, .
\end{align*}

\subsection{One-hot Encoding of Worlds}

Let us now generalize the one-hot encodings of propositional logic worlds to worlds of \firstOrderLogic{}.
To encode the boolean tensors $\worldindices$ describing \firstOrderLogic{}s as basis elements of a tensor space, we take the one-hot encoding
\begin{align*}
    \onehotmap \defcols \worldindexset \rightarrow \worldtensorspace
\end{align*}
defined by
\begin{align*}
    \onehotmapofat{\worldindices}{\worldvariables}
    =
    \left(\bigotimes_{\seccatenumeratorin}\bigotimes_{\indindexof{\seccatenumerator}\in[\inddim]^{\indorderof{\seccatenumerator}+1}}
        \onehotmapofat{\catindexof{\seccatenumerator,\indindexof{\seccatenumerator}}}{\catvariableof{\seccatenumerator,\indindexof{\seccatenumerator}}}\right)
    \otimes\left(\bigotimes_{\atomenumeratorin}\bigotimes_{\indindexof{\atomenumerator}\in[\inddim]^{\indorderof{\atomenumerator}}}
        \onehotmapofat{\catindexof{\catenumerator,\indindexof{\catenumerator}}}{\catvariableof{\catenumerator,\indindexof{\catenumerator}}} \right)
\end{align*}
This is a tensor of order $\catorder\cdot\inddim^{\indorder}$, in a tensor space of dimension $2^{\left(\catorder\cdot\inddim^{\indorder}\right)}$.
Storage of such tensors in naive formats would not be possible.
However, the basis $\cpformat$ format discussed in \charef{cha:sparseRepresentation} still provides storage with demand linear in the order $\catorder\cdot\inddim^{\indorder}$.


\subsection{Probability Distributions}

Having established the formalism of one-hot encodings also in the case of \firstOrderLogic{} worlds, we can now proceed with the definition of distributions and formulas, analogously to the development in \parref{par:one}.
Probability distributions over worlds coinciding on their domain are then non-negative and normalized tensors $\probat{\worldvariables}$ where each coordinate of a world $\worldindices$ is captured by a boolean random variable $\catvariableof{\atomenumerator,\shortindindices}$, indicating whether the $\atomenumerator$-th predicate holds on the object tuple indexed by $\shortindindices$.

% High-dimensional - watch out for repetitions!
We notice that by definition these probability distributions are distributions of
\begin{align*}
    \left(\sum_{\catenumeratorin}\inddim^{\indorderof{\catenumerator}}\right) +  \left(\sum_{\seccatenumeratorin}\inddim^{\indorderof{\seccatenumerator}}\right)
\end{align*}
Booleans.
% One-hot encodings minimal
Unfortunately, it is not possible to design encoding spaces of smaller dimension, when our aim is to get any distribution over possible worlds by an element in the encoding space.
This is due to the fact, that one-hot encodings provide a basis in the tensor space, as will be shown in \charef{cha:coordinateCalculus}.
The reason for the large encoding space dimension is therefore rooted in the equal number of possible worlds and not in an overhead in the dimension of the one-hot encoding space.
We will later in this chapter investigate methods to handle such high-dimensional distributions in the formalism of exponential families.



